\chapter*{Introducción}
\addcontentsline{toc}{chapter}{Introducción}
\markboth{Introducción}{Introducción}

El Método de los Elementos Finitos (MEF) es una de las herramientas centrales del análisis estructural y de la ingeniería computacional \autocite{zienkiewicz2013fem,cook2002concepts}, y su enseñanza forma parte de los planes de estudio de grado en ingeniería \autocite[p.~1160]{perezsantiago2023fem}. Quien lo aprende debe articular tres planos que suelen presentarse por separado: la formulación matemática continua (elasticidad lineal, ecuaciones de equilibrio); su discretización mediante elementos isoparamétricos e integración numérica; y la interpretación física de los resultados, es decir, de los campos de desplazamiento, deformación y tensión (esfuerzo). Que un elemento sea isoparamétrico significa que unas mismas funciones describen su forma y el desplazamiento en su interior. Y \emph{educativo} califica aquí el propósito de diseño del producto, no un efecto medido sobre quien lo usa: exponer el \emph{procedimiento de cálculo} ---la secuencia de etapas que va de las funciones de forma a la recuperación de tensiones--- además de entregar su resultado. Este trabajo desarrolla EduFEM, un software educativo de escritorio que reúne esos tres planos en un único entorno interactivo para problemas bidimensionales de elasticidad lineal estática, con elementos cuadriláteros de cuatro nodos (Q4) y de nueve nodos (Q9).

\section*{Situación problemática}
\phantomsection\addcontentsline{toc}{section}{Situación problemática}

La enseñanza del MEF tropieza con dos dificultades que se refuerzan entre sí. Por un lado, los textos clásicos del método, como los de Bathe y Reddy \autocite{bathe2014fem,reddy2006introduction}, desarrollan la teoría con un grado de abstracción considerable. Las funciones de forma, la matriz Jacobiana, la matriz de deformación-desplazamiento y la integración de Gauss aparecen como pasos de una sucesión algebraica, y su conexión con la geometría concreta de un elemento no siempre es evidente para quien se inicia. Por otro lado, el software comercial de análisis opera como una caja negra: el usuario define geometría, material y cargas y obtiene resultados, pero el procesamiento intermedio que transforma esas entradas en tensiones permanece oculto \autocite[p.~1160]{perezsantiago2023fem}. El antecedente educativo más directo de este trabajo lo describía ya en 1998 en los mismos términos: con los programas comerciales, el estudiante queda limitado a introducir los datos del problema e interpretar los resultados, sin acceso a los detalles del análisis \autocite[p.~246]{suarez1998edelas2d}.

Esa opacidad tiene una consecuencia técnica concreta: un análisis cuyos resultados no pueden seguirse hasta sus datos ni comprobarse paso a paso. En la práctica profesional, un cálculo estructural se entrega con su memoria de cálculo precisamente para que un revisor pueda seguirlo y comprobarlo. El software comercial, en cambio, entrega los resultados del análisis por elementos finitos sin la memoria de su procedimiento interno: mirando el programa no es posible saber de dónde sale una tensión nodal, ni si las tensiones de los nodos cierran con las calculadas en el interior de los elementos. La pertinencia del problema para la formación está documentada. Un estudio de consenso entre expertos de la industria y de la academia concluye que el egresado debe saber planificar, verificar y validar sus análisis, no solo ejecutarlos \autocite[pp.~1159, 1171]{perezsantiago2023fem}.

Los recursos educativos existentes no cierran esa brecha para la elasticidad plana. Entre los recursos revisados en la \autoref{sec:antecedentes}, ninguno documenta a la vez los ocho atributos que allí se enumeran, y la \autoref{tab:comparativa} compara esos recursos atributo por atributo. El vacío se concentra en dos rasgos que ninguno de los recursos didácticos revisados documenta: una memoria de cálculo reproducible a mano y una verificación y validación publicada junto con la herramienta. El antecedente más cercano, que sí expone el procedimiento de cálculo completo sobre el modelo del usuario, muestra en sus capturas una interfaz en inglés. A ello se suma que las herramientas comerciales no están diseñadas con un propósito didáctico explícito \autocite[p.~1160]{perezsantiago2023fem} y que su uso en el aula en idioma español es limitado. Esa es la situación de la que parte el trabajo: en la investigación tecnológica el problema no se elige libremente, sino que resulta de interpretar y diagnosticar una realidad concreta \autocite[p.~83]{garciacordoba2005tecnologica}. La insuficiencia de una herramienta que exponga el procedimiento de cálculo de modo que el análisis pueda seguirse hasta sus datos y comprobarse paso a paso es la contradicción que origina el trabajo \autocite[pp.~7, 14]{alvarez_metodologia}.

\section*{Formulación del problema}
\phantomsection\addcontentsline{toc}{section}{Formulación del problema}

¿De qué manera la escasa transparencia del procedimiento de cálculo del software de análisis por elementos finitos podrá incidir en la baja trazabilidad y verificabilidad de ese procedimiento, en la gestión 2026?

\section*{Objeto de estudio y campo de acción}
\phantomsection\addcontentsline{toc}{section}{Objeto de estudio y campo de acción}

El \textbf{objeto de estudio} \autocite[p.~8]{alvarez_metodologia} es el análisis de medios continuos en elasticidad lineal plana por el Método de los Elementos Finitos: el procedimiento de cálculo que va del mapeo isoparamétrico a la recuperación de tensiones con elementos cuadriláteros Q4 y Q9, sobre el que se construye el motor, se miden las variables y se contrastan los resultados. El \textbf{campo de acción}, la parte de la situación problemática que el trabajo se propone mejorar, son los medios de cálculo con que ese análisis se enseña: el software educativo que expone el procedimiento y la memoria de cálculo que lo documenta sobre el modelo del usuario.

\section*{Delimitación}
\phantomsection\addcontentsline{toc}{section}{Delimitación}

El problema se delimita en cuatro planos:
\begin{itemize}
    \item en lo disciplinar, a la elasticidad lineal plana con elementos cuadriláteros Q4 y Q9 (véase \emph{Alcance y limitaciones});
    \item en lo institucional, a la formación de grado en ingeniería civil, con la Carrera de Ingeniería Civil de la Universidad Autónoma Tomás Frías como destinataria inmediata de la herramienta;
    \item en lo espacial, al aula y a la computadora personal del estudiante, sin requerir laboratorio de cómputo ni licencias;
    \item en lo temporal, al desarrollo, la verificación y la validación del software realizados durante la gestión 2026.
\end{itemize}

\section*{Justificación}
\phantomsection\addcontentsline{toc}{section}{Justificación}

El desarrollo de EduFEM se justifica en cuatro planos complementarios. En el plano académico, la herramienta acorta la distancia entre la formulación teórica del MEF y su ejecución numérica: deja a la vista, sobre el elemento real que el usuario construyó, los pasos intermedios que el software profesional encapsula. Dos trabajos sobre enseñanza del método reportan resultados favorables de la simulación interactiva y de la construcción paso a paso de las herramientas de cálculo \autocites[p.~168]{lee2015interactive}[p.~1023]{bishay2020teaching}, con las salvedades que examina la \autoref{sec:fundamentos-pedagogicos}. Esa literatura orienta las decisiones de diseño de EduFEM.

En el plano práctico, EduFEM está concebido como apoyo a la cátedra. Es gratuito y de código abierto —su código fuente está disponible públicamente en línea\footnote{Repositorio del proyecto: \url{https://github.com/Sucullani/SoftwareED}}—, está íntegramente en español y organiza su interfaz según las tres fases habituales del análisis (pre-proceso, proceso y post-proceso), las mismas del software profesional. Además genera automáticamente una memoria de cálculo: un documento que reconstruye el procedimiento con sustitución numérica paso a paso, sobre los valores del modelo del propio usuario. Esta capacidad no se documenta en los antecedentes de software educativo revisados (\autoref{sec:antecedentes}). Con ella, el estudiante puede rehacer a mano, o en una hoja de cálculo, las operaciones que la memoria desarrolla ---de las funciones de forma a la rigidez elemental, su ensamblaje y la recuperación de tensiones--- y contrastarlas con las del software (\autoref{sec:memoria-io}).

En el plano técnico, construir un software propio en lugar de adoptar uno existente responde a la necesidad de controlar qué se expone del motor de cálculo. Un motor desarrollado sobre bibliotecas numéricas de código abierto \autocite{harris2020numpy,virtanen2020scipy} permite exhibir las matrices y las operaciones intermedias sin las restricciones de un producto cerrado, y verificar su exactitud frente a referencias reconocidas. La elección del lenguaje Python obedece a tres razones: no exige licencia al estudiante; reúne en un solo lenguaje el cálculo numérico, la interfaz gráfica y la generación de documentos; y permite distribuir el software como un instalador autónomo, que no requiere tener Python instalado.

En el plano profesional, el trabajo ofrece un ejemplo documentado de verificación y validación de un programa de cálculo estructural, el procedimiento del que depende la confianza en los resultados de cualquier programa de esa clase.

\section*{Objetivos}
\phantomsection\addcontentsline{toc}{section}{Objetivos}

\subsection*{Objetivo general}
\phantomsection\addcontentsline{toc}{subsection}{Objetivo general}

Desarrollar un software educativo de elementos finitos para el análisis estructural, empleando el lenguaje de programación Python, debido a la escasa transparencia del procedimiento de cálculo del software de análisis, para mejorar la trazabilidad y la verificabilidad del procedimiento de cálculo.

\subsection*{Objetivos específicos}
\phantomsection\addcontentsline{toc}{subsection}{Objetivos específicos}

\begin{enumerate}
    \item Fundamentar teóricamente el MEF en elasticidad plana con elementos Q4 y Q9, y el estado del arte del software para su enseñanza, para fijar los requisitos de la herramienta.
    \item Desarrollar EduFEM en Python ---motor de cálculo, interfaz, módulos educativos y memoria de cálculo--- para exponer el procedimiento de cálculo etapa por etapa sobre el modelo del usuario.
    \item Verificar y validar EduFEM, midiendo la cobertura de su procedimiento de cálculo y la exactitud de su motor frente a soluciones exactas y referencias externas, para establecer si ese procedimiento es trazable y verificable.
\end{enumerate}

\section*{Hipótesis}
\phantomsection\addcontentsline{toc}{section}{Hipótesis}

El desarrollo de EduFEM permitirá mejorar la transparencia del procedimiento de cálculo por elementos finitos ---de caja negra a procedimiento a la vista---, elevando su trazabilidad y su verificabilidad.

\section*{Variables}
\phantomsection\addcontentsline{toc}{section}{Variables}

La hipótesis relaciona tres variables, que la \autoref{fig:variables} representa:
\begin{itemize}
    \item \textbf{Variable independiente (causa):} la transparencia del procedimiento de cálculo, es decir, la exposición que el software hace de su procedimiento de cálculo interno, con sus magnitudes intermedias. Tiene dos niveles: \emph{caja negra}, que solo muestra datos y resultados, y \emph{procedimiento a la vista}, etapa por etapa.
    \item \textbf{Variable dependiente (efecto):} la trazabilidad y la verificabilidad del procedimiento de cálculo. La \emph{trazabilidad} es la propiedad por la cual cada resultado puede seguirse hacia atrás, etapa por etapa, hasta los datos del modelo, con todas las magnitudes intermedias a la vista. La \emph{verificabilidad} es la propiedad por la cual cada magnitud, intermedia o final, puede comprobarse contra un patrón independiente ---un cálculo manual, una solución analítica u otro programa--- y esa comprobación da conforme dentro de una tolerancia fijada de antemano. Son las dos propiedades que, en la práctica profesional, la memoria de cálculo busca asegurar.
    \item \textbf{Variable interviniente (solución):} EduFEM, la solución tentativa con que se actúa sobre la variable independiente.
\end{itemize}

Las dos primeras son atributos del procedimiento de cálculo, que es la unidad de análisis: se observan en el software y en los casos de estudio, no en personas. La relación entre ellas no es una tautología: exponer el procedimiento no asegura que lo expuesto esté completo ni que sea correcto, como muestra el defecto que este mismo trabajo detectó y corrigió en la extrapolación de tensiones a los nodos (\autoref{sec:leccion-extrapolacion}). La \autoref{sec:variables-met} operacionaliza las variables, y la \autoref{sec:validacion-resultados} fija de antemano los criterios con que se comprueba la hipótesis en sus dos cláusulas, una por dimensión de la variable dependiente.

\section*{Preguntas de investigación}
\phantomsection\addcontentsline{toc}{section}{Preguntas de investigación}

De la hipótesis se desprenden dos preguntas de investigación, una por dimensión de la variable dependiente:
\begin{itemize}
    \item \textbf{PI-1} (trazabilidad). ¿Qué organización del software permite seguir cada resultado hasta sus datos, etapa por etapa, sobre el modelo del usuario, y cubrir el contenido que exige un consenso publicado de expertos?
    \item \textbf{PI-2} (verificabilidad). ¿Reproduce el motor, dentro de tolerancias fijadas de antemano, las soluciones exactas, las referencias externas y los fenómenos numéricos que predice la teoría?
\end{itemize}

\section*{Metodología de la investigación}
\phantomsection\addcontentsline{toc}{section}{Metodología de la investigación}

Por su finalidad, este trabajo es \emph{propositivo}; por su naturaleza, es una \emph{investigación tecnológica} \autocite[p.~91]{garciacordoba2005tecnologica}. El método es el de la \emph{ciencia del diseño}, en la que el conocimiento sobre un problema y su solución se obtiene construyendo el producto tecnológico y evaluándolo con rigor \autocite[pp.~75, 83]{hevner2004design}. El enfoque es \emph{cuantitativo} y el nivel, \emph{descriptivo y comparativo}; la \autoref{sec:proceso-met} desarrolla esa clasificación, el método y sus seis actividades. Por tratarse de una investigación tecnológica, la hipótesis es la solución tentativa a un problema concreto, y se aprueba o no en la propia práctica de construir esa solución y ponerla a prueba \autocite[p.~85]{garciacordoba2005tecnologica}.

Este trabajo evalúa el software: su diseño, la exposición de su procedimiento de cálculo y la exactitud de sus resultados. Sus variables se miden sobre el software y sobre los casos de estudio. El diseño metodológico completo se desarrolla en la \autoref{sec:metodologia}.

\section*{Alcance y limitaciones}
\phantomsection\addcontentsline{toc}{section}{Alcance y limitaciones}

El alcance de EduFEM está delimitado explícitamente para lograr profundidad sobre un dominio acotado. En este documento, la expresión \emph{análisis estructural} se emplea en su acepción restringida de análisis de tensiones y deformaciones de elementos estructurales modelados como medios continuos bidimensionales, mediante la teoría de la elasticidad lineal. No comprende el análisis de sistemas estructurales de barras —pórticos o cerchas—, ni la formulación de placas, cáscaras o problemas tridimensionales.

La herramienta resuelve problemas \emph{bidimensionales} de \emph{elasticidad lineal estática}, en las dos idealizaciones clásicas de la mecánica de sólidos \autocite[p.~11]{timoshenko1970elasticity}: la \emph{tensión plana}, apropiada para cuerpos delgados cargados en su plano, y la \emph{deformación plana}, apropiada para cuerpos prismáticos largos cuya sección y cuya carga no varían a lo largo del eje. Dentro de ese dominio quedan problemas planos habituales en la ingeniería civil: en tensión plana, las vigas de gran peralte y los muros cargados en su plano; en deformación plana, las secciones de presas, de muros de contención y de túneles, donde una rebanada representa todo el sólido.

La discretización se realiza con \emph{elementos cuadriláteros isoparamétricos} de cuatro nodos (Q4) y de nueve nodos (Q9), que cubren respectivamente la interpolación bilineal y la cuadrática completa y permiten contrastar el efecto del orden de interpolación sobre la precisión \autocites[pp.~137-138]{hughes2000fem}[pp.~164, 183-184]{onate2009structural}. La pareja Q4--Q9 es el par mínimo que contrasta dos órdenes de interpolación dentro de una sola familia de funciones de forma: el Q4 es bilineal y el Q9, bicuadrático. Ambos reproducen los polinomios lineales completos; la diferencia no está en la completitud lineal sino en el orden, que es lo que gobierna el comportamiento en flexión. Mantener una sola familia conserva además las matrices del seguimiento paso a paso en un tamaño legible. Los elementos triangulares y los cuadriláteros de ocho nodos quedan fuera de esta versión, y su incorporación se plantea en las recomendaciones. El producto es un \emph{software de escritorio} de carácter \emph{educativo}, desarrollado íntegramente en \emph{español}, con una interfaz diseñada bajo criterios de sencillez y claridad.

La elección del dominio bidimensional no es solo una restricción de alcance, sino una decisión de diseño. El continuo plano ocupa el punto intermedio entre el análisis unidimensional de barras y el análisis tridimensional de sólidos. El elemento de barra, con el que suele iniciarse la enseñanza del método, no alcanza para mostrar los conceptos centrales de la formulación isoparamétrica: el mapeo entre coordenadas naturales y físicas, la matriz Jacobiana, la matriz de deformación-desplazamiento y la cuadratura en dos direcciones solo aparecen al pasar al continuo. El análisis tridimensional sí contiene esos conceptos, pero el tamaño de sus matrices los oscurece: la matriz constitutiva crece a $6\times6$, la de deformación-desplazamiento a $6\times24$ en un hexaedro de ocho nodos y la rigidez elemental a $24\times24$ \autocite[p.~218]{cook2002concepts}. Esas dimensiones vuelven impracticable el seguimiento paso a paso que constituye el núcleo de la herramienta. El plano es la dimensión mínima en la que el Jacobiano del mapeo isoparamétrico ya es una matriz y la cuadratura recorre dos direcciones, y a la vez conserva matrices de tamaño legible —$\bm{D}$ de $3\times3$ y $\bm{B}$ de $3\times8$ en el elemento Q4—, lo que permite exhibir cada operación. El dominio plano actúa así como un puente: la formulación isoparamétrica, la cuadratura de Gauss y el ensamblaje que el software expone en dos dimensiones son los mismos que rigen el caso tridimensional \autocite{bathe2014fem,zienkiewicz2013fem}, planteado como línea de continuación del proyecto.

Las limitaciones fijadas de antemano son las siguientes:
\begin{itemize}
    \item El software no aborda análisis tridimensional, comportamiento no lineal (del material o geométrico), efectos dinámicos ni dependientes del tiempo.
    \item El catálogo de materiales se restringe al comportamiento elástico, isótropo y lineal, definido por el módulo de elasticidad y el coeficiente de Poisson.
    \item La malla se construye manualmente —por tablas, por dibujo en el lienzo o por importación de un archivo DXF— o por generación estructurada en los casos de estudio: el software no incorpora mallado automático.
    \item La validación se realiza contra soluciones analíticas, casos clásicos de la literatura y un software comercial, no contra ensayos físicos.
\end{itemize}
Estas fronteras responden a una decisión deliberada: privilegiar la transparencia y la solidez de un núcleo bien acotado por sobre la cobertura amplia de funcionalidades. Las limitaciones que el propio desarrollo puso de manifiesto se discuten en la \autoref{sec:interpretacion}.

\section*{Estructura del documento}
\phantomsection\addcontentsline{toc}{section}{Estructura del documento}

El documento se organiza en tres capítulos, uno por objetivo específico, además de esta introducción, de las conclusiones y de siete anexos. Cada capítulo lleva por título su objetivo. Los tres capítulos corresponden, en ese orden, al marco teórico, al diseño e implementación del modelo y a la presentación y el análisis de resultados que prevé el índice de la tesis de grado en la Carrera.

El \autoref{cap:marco_teorico} cumple el primer objetivo. Abre con los antecedentes y el estado del arte del software para la enseñanza del MEF, de los que salen los requisitos de la herramienta, y con los principios que orientan la exposición del procedimiento. Desarrolla luego los fundamentos del método aplicados a la elasticidad lineal bidimensional, la calidad de malla y los criterios de verificación y validación.

El \autoref{cap:diseno} cumple el segundo. Abre con el diseño metodológico, del tipo de investigación a los criterios de aceptación con que se contrasta la hipótesis, y describe después el software, de sus requisitos a la memoria de cálculo.

El \autoref{cap:resultados} cumple el tercero. Ejecuta la verificación y la validación numérica frente a soluciones analíticas y a un software comercial, mide la cobertura del procedimiento y del contenido, interpreta los hallazgos y contrasta la hipótesis.

Las conclusiones siguen la misma correspondencia: una conclusión y una recomendación por capítulo, además del contraste de la hipótesis. Los anexos reúnen:
\begin{itemize}
    \item el manual de instalación (\autoref{anexo:instalacion});
    \item el manual de uso (\autoref{anexo:manual});
    \item una selección de listados de código del motor (\autoref{anexo:codigo});
    \item los resultados extendidos de verificación y validación (\autoref{anexo:vyv});
    \item el modelo de validación en SAP2000 con la solución analítica de Timoshenko (\autoref{anexo:sap2000});
    \item los formatos de intercambio de datos (\autoref{anexo:formatos});
    \item una memoria de cálculo resuelta paso a paso sobre el ejemplo canónico (\autoref{anexo:memoria}).
\end{itemize}

Quien no maneje el método puede seguir el argumento sin el desarrollo matemático que va de la \autoref{sec:formulacion-debil} a la \autoref{sec:calidad-malla}, con la Nomenclatura a mano para los símbolos: el problema, el método de la investigación y los resultados se siguen sin él.

\section*{Normas de citación y presentación}
\phantomsection\addcontentsline{toc}{section}{Normas de citación y presentación}

Las citas y la lista de referencias siguen la norma Vancouver del Comité Internacional de Editores de Revistas Médicas (ICMJE) y de la Biblioteca Nacional de Medicina de los Estados Unidos (NLM) en su principio: numeración arábiga por orden de primera mención. De su forma se apartan en tres decisiones, que se declaran. El número va entre corchetes, y no en superíndice, para poder acompañarlo con la página citada; el título de la revista se escribe completo, y no abreviado según el catálogo de la NLM, para no obligar al lector a descifrar abreviaturas; y se conserva la conjunción antes del último autor. Los asientos siguen, además, la puntuación del estilo numérico del gestor de referencias (\texttt{biblatex}), no la de la NLM: el título del artículo va entre comillas y conserva sus mayúsculas, y la revista va precedida de «En:», con el volumen, el número y el año en la forma «31.5 (2023)». Dentro de un corchete, la coma separa referencias sin localizador, o una referencia de su página, y el punto y coma separa las referencias cuando alguna de ellas lleva localizador. El localizador de página acompaña a toda cita que sostiene una ecuación, un valor numérico o un umbral, y también a la que sostiene una definición o una atribución concreta; se omite en las citas de encuadre. La presentación del documento sigue la séptima edición de las normas de la Asociación Estadounidense de Psicología (APA), con tres excepciones: los capítulos y las secciones van numerados, porque el texto remite a ellos por su número; el papel es A4; y el texto va justificado, con partición silábica. La fuente es Times New Roman de 12 puntos, una de las que la norma admite.
